\documentclass[11pt]{article}
\usepackage{setspace}
\usepackage{natbib}
\bibpunct[: ]{(}{)}{;}{a}{}{,}
\usepackage[top=1.75in, left=1.3in, right=1.3in, bottom=1.15in]{geometry}

\title{Network Transmission versus Culture Conversion and the Dynamic Stability of Cultural Tastes}
\author{Omar Lizardo}
\date{\today}

\begin{document}
\maketitle

\newpage
\begin{abstract}
Network models of cultural transmission and acquisition assume that current tastes are largely a function of the individual's existing and previous network environment.  Cultural conversion models on the other hand rely on the proposition (taken from Bourdieu's ``cultural capital'' approach) that tastes are driven by durable dispositions acquired throughout childhood and into early to middle adolescence and that these tastes are mobilized to construct and select appropriate network relations over time.  Given the volatility of social relations uncovered in recent longitudinal research on personal networks, the extant network transmission model implies highly malleable or unstable tastes over time.  The cultural conversion model on the other hand, predicts little over-time volatility in cultural preferences.  In this paper I attempt to adjudicate the appropriateness of these models of the culture-network linkage by evaluating current research and theory on the subject of dynamic taste stability.  I conclude that there is very little support for the implication of dynamic taste volatility.  Instead, the evidence appears to suggest that cultural dispositions are relatively stable, while network relations are instead dynamically volatile, consistent with the culture-conversion model's implications.  
\end{abstract}•

\newpage
\onehalfspacing
\section{Introduction}
\label{sec:intro}

Prevailing models of taste formation, transmission and decay in the sociology of culture posit that current tastes are largely a product of the individual's current network configuration.  The standard assumption is that tastes are transmitted through existing network ties, in particular those social connections that are characterized by a high degree of sociodemographic similarity---``status homophily'' \citep{lazarsfeld_merton54}---between ego and alter \citep{mcpherson_etal01}.  From this perspective, birds of a feather not only flock together, but largely due to this mutual flocking they end up singing, going to the movies, and reading together \citep{mark98b, mark03}. A key explanatory strength of the network model is its ability to account for the fact that the distribution of tastes across persons in socio-demographic space is ``lumpy'' and clustered in proximal regions of this space.  The model explains this by suggesting that the process of cultural taste formation and transmission is best envisioned as a ``local bandwagon'' process of network-based imitation \citep{mark03}. Individuals are construed as highly susceptible to the influence of their current contacts and by implication tastes are seen not as enduring aspects of the person but as malleable, ``thin'' contents liable to change when those network ties in turn change in their turn.  

Very different assumptions are at core of the cultural capital model of taste formation and acquisition \citep{bourdieu84, holt98, lizardo06a}. Instead of positing tastes as highly malleable contents at the perennial mercy of a constantly shifting network configuration, cultural tastes are conceived as highly durable dispositions toward broad types of cultural objects and performances \citep{frith98}.  These dispositions are formed in the family and school environment throughout childhood and early to middle adolescence \citep{smith95}.\marginpar{add reference to Holbrook paper on music taste and age}  Once set, they are reasonably resistant to change.  The reason for this is that as \citep{holt97} notes following \citet{bourdieu86}, tastes are part the person's embodied cultural capital and thus represents a highly resilient aspect of the individual's relationship to cultural objects \citep{bourdieu84}, encoded in the memory traces, semi-conscious affective tendencies and habitual dispositions responsible for practical action \citep{bourdieu90a}.

\subsection{Outline of the argument}
In this paper{\ldots} 

\section{Network versus cultural capital theories on dynamic taste acquisition and change}
\label{sec:net-vs-cult}
\subsection{The Network Transmission Model}
\label{subsec:net}

Contemporary network theories of cultural taste and attitude formation and change (and by implication decay) suggest that the contemporary taste profile of an individual is largely a function of the tastes held by the other persons in her surrounding social network \citep{erickson82, erickson88}.  The key underlying assumption is that social ties serves as the ``conduits'' through which new cultural practices are transmitted and by way of which individuals learn these new cultural practices from their social contacts \citep{mark98a, carley91}.  This proposition is shared by both straightforward network accounts \citep[e.g.][]{cardon_granjon05, erickson96, kane04} and more ambitious ecological-network theories of cultural taste \citep{mark98a, mark03, mcpherson04}. \citet[323--324]{mark03}  for instance speaks of social ties to friends and family as the most important sources of cultural taste acquisition:  ``[p]eople's friends and family{\ldots}influence the leisure activities they participate in{\ldots}participation in various sports, visiting art museums, and gardening are{\ldots}leisure activities that people often lean from parents, siblings, and/or friends.''

Network theories of taste transmission implicitly presume a relatively stable ``social structure'' coupled to relatively unstable patterns of cultural taste, otherwise it would be hard to think how something that varies at a more rapid pace ``causes'' or ``affects'' something that is more stable than itself \citep{simon52}.  It bears mentioning that this relative stability assumption is shared by the entire family of network theories that conceive of social structure as the pipes through which contents, or information ``flows'' and diffuses from one actor to another as \citet{borgatti05} has recently noted.  This implicit (and crucial) assumption of the stability of social structure, when this term is defined---as it is usually done in network theory---as the ``concrete'' patterns of network relations that bind individuals to one another \citep{berkowitz82, berkowitz88, wellman83, wellman88, bearman97, breiger04}, can certainly be questioned in light of recent longitudinal research on constancy of network relations over time.  This research---see the 1997 and 2005 special issues of \textit{Social Networks} (volume 19, number 1 and volume 27 number 4) and \citet{burt00} for a thorough review)---shows that instead of being constant, network relations are extremely volatile, even when the time frame is constrained to relatively short periods of time, such as a year or less \citet{bidart_degenne05, suitor_etal97}. Thus, what is truly surprising about over-time network ``stability'' over time is ``how unstable intimacy is'' \citep[47]{wellman_etal97}.  

Accordingly, most longitudinal studies of personal networks reveal surprisingly high levels of instability in personal networks over time \citep{lang00}, with anywhere approximately about 30 to 60\% of the personal network experiencing turnover in about a year's time.  In his own research on bankers, \citet[11]{burt00} finds that ``three in four of the 22,709 colleague relations at risk of decay{\ldots}are gone next year.'' This means that old ties are constantly being deleted---what \citet{burt00} refers to as ``tie decay''---and new ones formed constantly throughout the life-course \citet{suitor_etal97, wellman_etal97}.  \citet[57]{suitor_keeton97}, report that more than half of the associates...named at T1 were not named again at either T2 or T3.''   \citet[14]{morgan_etal97} conclude that their results ``point to a considerable degree of instability in who is present in the network over time. In particular, the average overlap rate of 55\% means that just over half of those who were named in one interview or the other appeared in both. Similarly, interpreting the average phi value of 0.34 as a test-retest correlation indicates only a limited ability to predict who was in the network from one time to the next.''

These findings imply that for most individuals a recurrent process of new contact selection, extant relationship ``fine tuning'' and possible deletion of old social ties is constantly at work \citep{lazarsfeld_merton54}.  This represents an important---if ill understood---facet underlying the temporal evolution of a person's relational micro-environment.   This dynamic acquires renewed salience in light of the observation that recently formed network connections appear to exhibit what \citet{burt00}, drawing on \citet[148--150]{stinchcombe65}, refers to as a ``liability of newness.'' This means that dyadic network ties tend to display a high probability of dissolution at early stages of the relationship, with the probability of termination steeply declining as the relationship matures. This proposition implies that the ``tie decay function'' (that is the probability of a tie being dissolved) is non-linearly decreasing in relationship age, with a high probability for young relationships which decreases as these dyadic ties ripen into close friendships.  \citet{burt00, burt02} finds evidence consistent with this hypothesis for both strong ties and network bridges (weak ties), respectively.

It is true that network stability is much more substantial in the ``core'' personal network, with the bulk of the instability concentrated on the periphery\citep{bidart_lavenu05, suitor_etal97, wellman_etal97}.  However, most longitudinal network studies show that the stable ``core'' network tends to be composed of a small set of alters, primarily kin.  The periphery of the network includes most of the persons usually referred to as (non-kin) ``friends.''  Studies show that there is considerable turnover even among those ties usually considered ``close friends'' \citep{bidart_lavenu05}. It is possible to revise network theories of taste transmission to suggest that taste are mostly transmitted through the small, relatively stable core and not through the larger, unstable periphery.  However, most contemporary network theories of taste do not make this distinction, and assume that taste are as equally likely to be learned from family and friends \citep{mark03}.  It is clear from recent studies on personal communities that the ``core'' kin network while mobilized for many purposes, including emotional and financial support, is not used for companionship and sociability \citep{simmel49} in any meaningful sense.  Most research and theory on culture consumption communities shows that if tastes are exploited in interaction, this is done precisely with those contacts that are usually named as ``friends'' and which are seen outside the home and the neighborhood \citep{dimaggio87, fiske87, frith98}.  However, there is nothing inherent in the logic of the network theory of taste transmission that makes it incompatible with current data suggesting large rates of network turnover and the liability of newness of social ties.  

The empirical fact of network volatility coupled with the network assumption of taste as being largely a function of the cultural choices of the individual's current set of contacts, implies large-scale volatility of those very same cultural tastes.  Thus if we take the network perspective seriously, we should expect that as friends change, so do tastes, so we should expect small to moderate correlations between an individual's taste at time t and her taste at time $t+n$. The reason for this is that as \citet{mark03} notes a taste that is not ``reinforced'' by social network influence through interaction is likely to be ``forgotten'' (e.g. lost).   This is a key assumption of the network based model of taste acquisition.  Thus, if we retain the network assumption that individual tastes are largely a function of the tastes displayed by the persons current network alters, and if we accept the fact that these alters are constantly change over the adult life course, then we should expect cultural tastes to themselves exhibit high volatility over time.
\subsection{The Culture Conversion Model}
\label{subsec:cult}

The cultural capital approach to taste formation provides a very different picture of the stability of tastes than that which can be gleaned from contemporary network theory.  Instead of implying that tastes are unstable, ``shallow'' contents liable to be replaced at the same pace that network relations are replaced, the cultural capital approach suggest that taste are ``deep'' durable dispositions and trans-temporally stable forms of embodied cultural capital \citep{bourdieu86, holt98}.  It is the relative stability of tastes vis a vis network relations that allows them to function as resources to help sustain and maintain these network relations over time---primarily by serving as topical resources in conversational interaction rituals \citep{dimaggio87, lizardo06a} providing a common focus of attention and affective commitment \citep{collins04}.  This helps to shift interactions away from purely ``instrumental'' purposes of information gathering and rational goal attainment to sociable interaction for its own sake \citep{granovetter02, simmel49}. 

\citet{lizardo06a} uses this stability assumption to suggest that tastes can in fact be seen as probable (reciprocal) drivers of network connections rather than being seen as exclusively determined by those connections.  Using two-stage instrumental variables and matching methods on cross-sectional data, that study uncovers evidence consistent with this reciprocal effect. Engagement in a wide variety of cultural practices increases network size net of standard controls and after statistically correcting for simultaneity bias.  This is in agreement with dynamic theories of cultural and network evolution, which rather than postulating a one-way avenue of causation going from networks to cultural tastes propose a reciprocal (over-time) two-way casual interaction between the two \citep{carley91, mark98a}. However, in the aforementioned study the stability hypothesis remains an empirically unverified background assumption.  Furthermore to reliance on cross-sectional data, the casual effect of cultural taste on network maintenance or even the more widely accepted effect of network turnover on cultural taste loss, has not been decisively established.
 
Thus, according to the cultural-capital theory inspired  conversion model of cultural into social capital \citep{dimaggio_mohr85, dimaggio_useem78a, lizardo06a}, the primary ways in which tastes are used to sustain social networks, is by allowing the individual to connect to alters that are both relationally close---when specialist or ``niche'' cultural tastes are concerned---and relationally distant---when talking about aptitude for ``popular'' cultural forms.  From this perspective, culture consumption and the ability to decode and appropriate symbolic goods produced in the urban, folk and ``culture-industry'' realms \citep{crane93} are important resources that individuals mobilize in the construction and maintenance of active social relationships \citep{cardon_granjon05, dimaggio87}.  The culture conversion model can in this particular respect be thought of as an elaboration and extension (to the realm of cultural taste) of  the classical \citep[e.g.]{lazarsfeld_merton54} process model of friendship formation, with similarity based on cultural taste and practices \citep{dimaggio87, mark03} playing the role that ``value-homophily'' played in the original formulation.    In this respect,  network models of taste formation through social network ties \citep{mark98b} can be thought of as partially based on the Lazarsfeld-Merton formulation.  However, these models raise what was explicitly a secondary dynamic in the original scheme---i.e. mutual ``value-adjustment'' through social interaction of initially different value dispositions \citep[33]{lazarsfeld_merton54}---to being the only operative, and thus primary dynamic.  It is clear however that for \citet[22--29]{lazarsfeld_merton54} ``selective'' processes whereby contacts are sifted based on similarity of attitudinal dispositions similar to those emphasized in the culture-conversion account were considered primary.

There is plenty of evidence that shows that tastes do play the role accorded to them (vis a vis network ties) in the culture conversion model.  As \citet[302]{cardon_granjon05}  note in their recent study of the role of culture consumption in the formation of social connections among young adults in France, ``{\ldots}a large proportion of ordinary conversational interaction is generally devoted to comments{\ldots}on cultural practices and recreational consumption. These are recurrent topics, present in a wide range of social situations{\ldots}Many social relations are constructed, in different ways, around cultural and recreational activities.''  As \citep[453]{fine79} noted in a classic paper, ``culture produces considerable discussion, and thus {\ldots}[its] content {\ldots}will become a resource in conversational interaction.'' Accordingly, those cultural goods which find an audience, will be deployed as the basis of conversation and will read the ``fodder for least-common denominator talk.'' We should expect that the consumption of aesthetic goods infuses local conversation-based interaction rituals with a common focus of attention allowing persons to spend time engaged in (Simmelian) interaction for its own sake\citeyearpar[443]{dimaggio87}. 

For instance, Elizabeth Long, in her study of women's reading groups in Houston, Texas, finds that cultural knowledge and taste become the primary content of social interaction. Consistent with the claim that connects the consumption of aesthetic products and sociability, she finds that reading group members ``tend to press books into service for the meanings that they transmit and the conversations they generate'' \citeyearpar[73]{long03}.  The primary function of aesthetic consumption experiences is to serve as the non-instrumental content that justifies the existence of status-segregated conversational rituals: ``{\ldots}talk about entertainment-the discussions of movies, music, sports, actors and athletes…make up so much of casual chitchat.'' This becomes even more important ``{\ldots}if we include talk about things people consume or seek out-food, clothes, houses and gardens, restaurants scenery, vacations. In some social circles, this entertainment talk makes up most of the content of conversations; in all classes, it is important" \citeyearpar[175]{collins75}. The Simmelian language of “social circles” is not accidental, since as \cite[797]{kadushin66} noted in his classic study on the subject, Simmelian urban social circles are produced and reproduced through shared cultural tastes---e.g. concert going, visiting the museum, etc.---and other activities intrinsically connected to culture-mediated interaction. 

From this perspective of the culture-conversion model therefore, while the effect of social network change on cultural taste loss is not denied, whatever relational (or “network”) effects on taste exist, must operate in the background of high-levels of cross-temporal taste stability and the reciprocal action of cultural taste on social networks.

 
\section{Evaluating current knowledge on the dynamic stability of tastes}
In which direction does the extant evidence point?  While there has been little empirical research on the dynamic stability of tastes through time, there is good reason to suppose that they are more stable than the network-transmission model indirectly implies.  Consider the following well-established findings:

\begin{enumerate}
\item  Most studies of the role of early family and school experiences on cultural capital have shown strong effects of arts participation, education and after-school training during adolescence on adult tastes even after controlling for subsequent educational attainment \citep{kracman96}.\marginpar{Talk about DiMaggio's 1982 paper} Hagan's \citeyearpar{hagan91} classic study of the effect of subcultural preferences on adult stratification outcomes implicitly supports a stability assumption. Hagan finds that involvement in a ``party subculture'' in adolescence has deleterious effects on the career advances of working class youth, but actually \emph{helps} the attainment chances of middle class youth, once the negative effects of these set of subcultural preferences on school performance is adjusted for. Hagan's concludes that it is precisely because the cultural habits obtained during this episode of ``drift'' continue to operate in adulthood that we find this class-specific enabling effect.

\item Studies that look at the patterning of taste across age groups show that specific genre preferences (e.g. for musical styles or even particular songs) are acquired early in adolescence and do not change dramatically afterwards. \citet{smith95} using data from the 1993 General Social Survey, shows that musical tastes are developed early in young adulthood and appear to be fairly distinct across age-cohorts.\marginpar{add quote from Smith and Holbrook}  

\item The few studies that are able to speak to the stability issue, indeed find support for the cultural-capital implication.  \citep[178]{bennett_etal99}, using qualitative evidence from retrospective interview reports on musical preferences in Australia find that ``{\ldots}many of those who were in middle age had retained  the music preference of their late teens and early adulthood.'' They conclude that while they ``{\ldots}cannot discount the possibility that some people will undergo radical changes in their tastes as they age, on bzalance the evidence appears to favor the explanation that the taste acquired in the formative years of a cohort persist into later life'' \citep[180]{bennett_etal99}.  In a recent study in Germany \citet{georg04}  reinterviewed a sample of respondents in 2002 which had previously been interviewed in 1983.  In the initial study, which was conducted when most respondents were adolescents, respondents provided information on a series of items regarding the participants preferences for various form of high-status cultural activities (e.g. attending a classical music or jazz concert; reading serious literature, etc.), most of which were then replicated in 2002. Drawing on the cultural-capital expectation of taste stability, Georg hypothesized that cultural preferences for high-status culture should be relatively stable throughout the life course. Using cross-lagged structural equations models he finds striking support this hypothesis ($\beta=0.61$, with $\beta=1.0$ being perfect stability and $\beta=0$ being no-stability).
\end{enumerate}•

Taken together, these findings suggest that while the effect of social network change on cultural taste loss is not denied, whatever relational (or ``network'') effects on taste exist, must operate in the background of high-levels of cross-temporal taste stability and the reciprocal action of cultural taste on social networks. Thus, contrary to the network transmission account, and in accordance with the imagery proposed in the culture-conversion model, tastes (in contrast to network ties) do not exhibit over-time volatility, but instead show high inter-temporal correlation and ``stickiness'' in the following section I explore some important implications of this conclusion. 

\section{Empirical and theoretical implications}
\subsection{Choice-homophily of network contacts based on pre-existing cultural dispositions}
Standard network theories of taste transmission draw on the notion of ``homophily'' (i.e. the empirical generalization that close friends tend to be similar in sociodemographic characteristics) to explain the fact that cultural tastes and practices tend to cluster in certain areas of social space \citep{mark03}.  From the standard point of view homogeneous social connections based on socio-demographic similarity pre-exist individual patterns of taste and this explains why cultural practices ``lump'' in certain areas of social space.  That is, the logic of network-based theories leads to them to attempt to explain the fact that cultural practices occupy delimited niches in social space by pointing to two interrelated processes:  (1) tastes tend to be associated with social background and (2) individuals tend to associate with those who are socio-demographically similar.  

Yet, there is sufficient qualitative \citep{cardon_granjon05, long03, wali_etal02} and quantitative \citep{lizardo06a, vaisey_lizardo10, wimmer_lewis10} evidence to suggest that just in the very same that tastes and subcultural preferences and values are transmitted through pre-existing network ties, new social ties are formed on the basis of pre-existing cultural dispositions \citep{lazarsfeld_merton54}.  This is the obverse of the traditional network-theoretic emphasis on the formation of new cultural tastes on the basis of preexisting social ties.  In fact, given the relative stability of cultural tastes in comparison to network connections, it is likely that this last process may in fact occur more often that the former process \citet{lizardo06a}.  In this respect (and in contradiction to standard network accounts), the ``lumpiness'' of cultural tastes in social space may be actually instead be the \emph{result} of an active---but not necessarily conscious---process of contact selection and relationship fine-tuning keyed around cultural compatibility \citep{lazarsfeld_merton54, vaisey_lizardo10, dimaggio93b}.

 This proposition is concordant with Burt's  observation of the ``liability of newness'' of recently formed ties, and with recent studies of changes in social networks of young adolescents and young adults \citep{bidart_degenne05, bidart_lavenu05, cardon_granjon05}.  If an initial ``feeling-out'' process (keyed to cultural compatibility) is characteristic of newly formed relationships, and if culturally incompatible contacts are likely to drop the relationship at early signs of incompatibility, then the higher mortality hazard of new relationships is readily explicable \citep{vaisey_lizardo10}.  Notice that if this is correct, then sociodemographic homophily could in the majority of cases be a by-product of cultural compatibility and not the other way around (this cultural compatibility may of course be itself due to socio-demographic of the family of origin [Bourdieu 1984; Illouz 1998]), as is strongly suggested in recent study of the cultural-underpinnings of racial homophily using dynamic network data \citep{wimmer_lewis10} and of the changing composition of adolescents networks using dynamic panel data \citep{vaisey_lizardo10}.  This is an insight that was always implicit in the theoretical logic of dynamic simulations of network formation based on shared cultural knowledge \citep{carley91, dimaggio93b}, but which has been occluded in recent theory and research due to the one sided emphasis put on the socio-structural determination of shared tastes.  
 
\subsection{Class-based heterogeneity in contact choice and structural determination of tastes}
Another important implication of the above considerations is the fact that the extent to which contacts are subject to a process of active selection based on taste itself varies according to the individual's command of cultural, educational and socio-economic resources.  Put simply, high cultural capital individuals, in particular those employed in prestigious occupations that demand of extensive degrees of cognitive and cultural skill not only have larger and sparser social networks \citep{lizardo06a}, but the proportion of their networks that is composed of constantly changing, loosely bounded social ties is larger and more heterogeneous than that of culturally disadvantaged individuals employed in routine, low-complexity occupations \citep{coser91, fischer82}.   In this respect, the endogenous effect of cultural capital in both the ability to select among different social ties and the ability to reach across the social structure to form those ties, is itself partially determined by position in social space. 
 
This is important, because the dominant way that ``social space'' has been conceptualized by network and organizational theorists \citep{hannan_etal03, mcpherson04} in relation to culture has been---following \citet{blau77a}---as a multidimensional manifold composed of variables such as age and occupational prestige.  The sociodemographic niche for a particular cultural or social form can then be given a metric definition in terms of these variables.  Social distance among two points in these space is conceived as being proportional to sociometric distance in social networks, that as being a proxy for the probability that two individuals will interact in real life \citep{mark98a}.  This presumes that the socio-demographic covariates that define social space are strictly exogenous (i.e. they do not interact in the statistical sense) with the propensity to deploy cultural knowledge reaped from consuming the cultural form in question in interaction to form new social relationships with persons located at relatively distant locations in social space.   Traditional network models also assume that there is no interaction between privilege locations in social space and the relative propensity to acquire or learn about specific cultural forms from proximate social contacts.  This process is seen as being a generic dynamic applicable to all individuals regardless of their socio-demographic profile.  

From the point of view of the culture-conversion model on the other hand the axes of social space as defined above are clearly composed of variables some of which (i.e. educational attainment) are themselves implicated in the extent to which tastes may be driven by social network dynamics.  Thus, from this point of view, occupying a privileged position in social space plays an important role in the degree to which individuals are able to loosen the limiting effects of network-based social constraint. This happens due the greater ability of those who posses higher levels of educational and socio-economic resources to deploy a more heterogeneous repertoire of cultural knowledge in everyday interaction.   This allows them to form and sustain social relationships that cut-across social dimensions of stratification and differentiation \citep{dimaggio87}, as well permitting them to create and maintain intimate ties with those who share their knowledge and tastes.  For this reason, the dimensions of ``social space'' cannot be considered strictly exogenous to the cultural contents whose sociodemographic niche they define.  Instead differential command of these cultural contents on the part of individuals may themselves ``warp'' these dimensions, such that individuals located at some privileged locations in this manifold find it easier to traverse a large distance in that space and thus connect to a wide range of others \citep{erickson96, lin_dumin86}. Individuals located in less privileged areas of the space, on the other hand, are socially and culturally distant from most individuals and cannot connect very well across their locally bounded social circles \citep{bryson97}. 

There is plenty evidence that is consistent with these considerations.  For instance \citet{wellman79, wellman_etal88, wellman_wortley90} and \citet{fischer82} show that as socio-economic and educational status increases, individuals are much more likely to have role-segmented, specialized relationships in which a clear differentiation is made between those ties designed for sociability and companionship (usually labeled as ``friends'') and social ties that are the result of non-negotiable statuses such as kinship ties.  These lines of research on personal community networks also show that individuals are much more likely to prefer to obtain companionship from voluntary rather than involuntary ties.  These are the social contacts that the culture conversion model predicts are likely to be formed and sustained through pre-existing forms of cultural compatibility.  These are also the social relationships that are likely to be sustained through interaction rituals in which the symbolic goods coming from artworlds and cultural industries figure as the primary subject of conversation \citep{cardon_granjon05, dimaggio87}.  Finally, these are also the same social ties that are continually undergoing high rates of turnover and being constantly recycled and reformed, and which thus require the recurrent deployment of cultural capital in interaction.  This goes a long way toward explaining the almost insatiable demand for novel symbolic goods among high cultural capital class fractions \citep{douglas_isherwood96}

In essence, for a relatively privileged cultural and economic elites social relationships have come to be partially disembedded from place and kin constraints \citep{dimaggio_mohr85, holt98}, and have taken the form of increasingly longitudinally unstable and thus constantly changing personal community networks\citep{wellman_etal88}.  This increases the chances that differential command of the ability to consume a wide variety of cultural forms becomes a driver of the recurrent formation of new network relations.  This is less likely to be the case for those  individuals who are restricted to a more homogeneous, kin-based, and locally constrained set of core alters who change relatively little over time.  Consistent with this claim, \citet{lizardo06a} shows that after holding cultural capital constant, the positive effect of education and socio-economic status on social network size is reduced significantly.  This has the curious implication that traditional network theory may more accurately describe the process of cultural taste formation among low cultural capital individuals subject to high levels of network constraint \citep{burt05} while the process of culture conversion may be a more accurate representation of the deployment and operation of cultural taste among culturally advantaged elites.  Consistent with this suggestions, a recent study by Lizardo \citeyearpar{lizardo11}, shows that the extent to which individuals specialize on a given cultural pursuit is associated with the extent to which they are likely to have contacts in their ``core discussion network'' who are also acquainted with one another.


\section{Concluding remarks}
This paper has attempted to advance current theory and research on the culture network linkage by evaluating our current  pertains directly to the empirical adequacy of different assumptions regarding the question of the stability of cultural tastes over time.  The question of taste stability is important, since it bears on the direct empirical implications of the two major extant frameworks purporting to shed light on the relationship between social ties and cultural resources:  network theory whether in its relational \citep{erickson96} or ecological \citep{mark98a, mark98b, mcpherson04} variants, and cultural capital theory \citep{bourdieu84, bourdieu86} (of which the culture-conversion model is a specific variant).  I argued that the current empirical record pointing to the accuracy of the cultural capital theory's assumption of stability and by implication of the misleading nature of the network-theoretic implication of instability has significant implications for how we theorize the statics and dynamics of cultural taste formation, transmission and decay.  Rather than the generic idea found in the traditional network model of the individual who is influenced by the tastes of those persons that he or she is connected to, the cultural capital derived hypotheses of relative taste stability and use of cultural taste for the formation of new relationships provides us with novel ways of conceptualizing long-standing issues in network theory and research (e.g. the dynamic origins of homiphily) as well as providing conceptual tools with which to theorize the inter-connection between culture, networks and social stratification in contemporary societies.
\newpage
\singlespacing
\footnotesize
\bibliographystyle{asr}
\bibliography{N:/Private/bib/main}

\end{document}





